\chapter{Weka GUI e Feature Selection}

\section{Introduzione a Weka}

\textbf{Weka} è un toolkit per il \textbf{machine learning} e il \textbf{data mining} scritto in Java e distribuito con licenza GNU. Viene usato per ricerca, didattica e applicazioni pratiche, ed è collegato al libro \emph{Data Mining} di Witten e Frank.

Le componenti principali di Weka sono:

\begin{itemize}
    \item \textbf{Explorer}: permette di esplorare visivamente i dati, applicare filtri e usare algoritmi di classificazione;
    \item \textbf{Experimenter}: consente di confrontare più algoritmi su uno o più dataset, anche con test statistici;
    \item \textbf{Knowledge Flow}: permette di progettare graficamente pipeline di machine learning;
    \item \textbf{Simple CLI}: interfaccia da riga di comando per eseguire direttamente classi Java di Weka.
\end{itemize}

\begin{notaprof}
Nel corso si usano soprattutto \textbf{Explorer} ed \textbf{Experimenter}. Explorer è utile per prendere confidenza con dataset, distribuzioni, filtri e classificatori; Experimenter è più adatto per confrontare più modelli in modo sistematico e ottenere risultati statisticamente confrontabili.
\end{notaprof}

\section{Formato dati in Weka}

Weka lavora principalmente con file \textbf{flat}, cioè dataset organizzati come tabelle relazionali singole.

\begin{notaslide}
I formati supportati includono \textbf{ARFF}, \textbf{CSV}, \textbf{C4.5} e \textbf{binary}.
\end{notaslide}

Il formato proprietario di Weka è \textbf{ARFF}, che si basa su tre parti principali:

\begin{itemize}
    \item \texttt{@relation}: nome del dataset;
    \item \texttt{@attribute}: definizione degli attributi, cioè nome e tipo di ogni colonna;
    \item \texttt{@data}: valori delle singole istanze, separati da virgola.
\end{itemize}

Weka distingue tra attributi \textbf{numerici} e attributi \textbf{nominali}. Per impostazione predefinita, l'ultimo attributo viene considerato la \textbf{classe target} da predire.

\begin{notaprof}
Nel progetto è consigliabile importare i CSV e salvarli subito in formato ARFF. Il CSV non specifica esplicitamente i tipi delle colonne, quindi Weka potrebbe interpretare alcuni attributi in modo errato. Inoltre, Weka non gestisce direttamente le scale ordinali, come \emph{low}, \emph{medium}, \emph{high}: queste vanno trasformate con attenzione, ad esempio tramite codifica binaria o numerica.
\end{notaprof}

\section{Explorer: preprocessing e classificazione}

Nel tab \textbf{Preprocess} dell'Explorer è possibile importare i dati tramite il pulsante \emph{Open file}. Dopo l'importazione, Weka permette di visualizzare gli attributi, le distribuzioni, il numero di valori distinti, i valori mancanti e alcune statistiche descrittive.

I filtri di preprocessing permettono di trasformare il dataset prima dell'addestramento.

\begin{notaslide}
Tra i filtri disponibili nelle slide sono citati \textbf{discretization}, \textbf{normalization}, \textbf{resampling} e \textbf{attribute selection}.
\end{notaslide}

Selezionando un attributo, Weka mostra informazioni come minimo, massimo, media, deviazione standard e un istogramma colorato rispetto alla classe target.

Nel tab \textbf{Classify}, Weka permette di scegliere diversi classificatori.

\begin{notaslide}
Tra i classificatori citati nelle slide compaiono alberi decisionali come \textbf{J48} e \textbf{RandomForest}, classificatori instance-based come \textbf{IBk}, SVM, logistic regression e Bayes nets, tra cui \textbf{NaiveBayes}.
\end{notaslide}

Weka include anche \textbf{meta-classifiers}, cioè classificatori costruiti sopra altri classificatori.

\begin{notaslide}
Tra i meta-classifiers sono citati \textbf{bagging}, \textbf{boosting} e \textbf{stacking}.
\end{notaslide}

Cliccando sul nome del classificatore è possibile aprire l'\emph{object editor} e configurare i parametri. Le principali opzioni di test sono:

\begin{itemize}
    \item \textbf{Use training set};
    \item \textbf{Supplied test set};
    \item \textbf{Cross-validation}, ad esempio 10-fold;
    \item \textbf{Percentage split}.
\end{itemize}

\section{Experimenter}

L'\textbf{Experimenter} permette di confrontare diversi learning schemes in modo più sistematico rispetto all'Explorer. È utilizzabile sia per problemi di classificazione sia per problemi di regressione.

\begin{notaprof}
L'Experimenter è particolarmente utile quando bisogna confrontare più classificatori, più dataset o più configurazioni. I risultati possono essere esportati, ad esempio in formato CSV, per analisi successive.
\end{notaprof}

\begin{notaslide}
Le slide indicano che l'Experimenter supporta \textbf{cross-validation}, \textbf{learning curve} e \textbf{hold-out}. Inoltre può iterare su diversi parametri e include strumenti di \textbf{significance testing}.
\end{notaslide}

Il significance testing permette di capire se la differenza tra due classificatori è statisticamente significativa, invece di limitarsi a confrontare valori medi di accuratezza.

\section{Feature Selection: obiettivi}

La \textbf{feature selection}, o selezione degli attributi, serve a scegliere un sottoinsieme di feature informative rispetto al task predittivo.

Gli obiettivi principali sono:

\begin{itemize}
    \item ridurre il costo di learning, diminuendo il numero di attributi;
    \item migliorare le performance rispetto all'uso del \emph{full attribute set};
    \item ridurre il costo di misurazione delle feature.
\end{itemize}

\begin{notaesame}
Nel contesto della Software Engineering, il costo più importante non è sempre il tempo di addestramento del modello, ma il costo di misurazione. Estrarre decine o centinaia di metriche per ogni classe può richiedere molto effort. Ridurre il numero di feature significa ridurre anche il costo di raccolta dei dati.
\end{notaesame}

Nei dataset di defect prediction è frequente avere metriche ridondanti o sovrapposte. Per esempio, LOC, numero di metodi e numero di attributi possono misurare indirettamente la dimensione della classe. La feature selection permette di mantenere solo le feature più utili.

\section{Filter approach}

Il \textbf{filter approach} valuta gli attributi partendo dai dati, senza considerare il classificatore che verrà poi usato.

I filtri assegnano alle feature un \textbf{merit score}, cioè un punteggio che misura quanto una feature sia utile per separare le classi o per descrivere la variabile target.

La valutazione può avvenire in due modi:

\begin{itemize}
    \item valutando ogni feature singolarmente;
    \item valutando sottoinsiemi di feature.
\end{itemize}

Valutare le feature singolarmente è veloce, ma può essere sub-ottimale perché ignora le relazioni tra attributi. Valutare i subset è più potente, ma introduce un problema combinatorio: il numero di possibili sottoinsiemi cresce in modo esponenziale.

\begin{notaprof}
I filtri sono molto veloci e utili quando il numero di attributi è grande. Il limite è che il punteggio prodotto, ad esempio un valore di Info Gain, non sempre indica chiaramente dove fissare la soglia per eliminare le feature.
\end{notaprof}

\section{Info Gain e Spearman}

\subsection{Information Gain}

L'\textbf{Information Gain} misura quanta informazione una feature fornisce sulla classe target. Più precisamente, quantifica la riduzione dell'incertezza, misurata tramite entropia, quando si conosce il valore di un attributo.

Un valore alto di Information Gain indica che la feature è utile per separare le istanze in classi distinte.

\begin{notaprof}
Una feature con alto Information Gain aiuta il classificatore perché riduce l'incertezza su quale classe assegnare a un'istanza.
\end{notaprof}

\subsection{Spearman correlation}

La \textbf{Spearman correlation}, o coefficiente rho di Spearman, è una misura non parametrica della forza e della direzione dell'associazione tra due variabili ordinate.

\begin{notaslide}
A differenza di Pearson, che misura relazioni lineari sui valori grezzi, Spearman valuta relazioni monotone sui ranghi. Il coefficiente varia da \(-1\) a \(+1\).
\end{notaslide}

Una relazione monotona esiste quando una variabile tende ad aumentare al crescere dell'altra, oppure tende a diminuire al crescere dell'altra, anche se non necessariamente a ritmo costante.

\begin{notaprof}
Metriche come Spearman sono interpretabili, ma non risolvono completamente il problema della ridondanza. Due feature possono essere entrambe molto correlate con la classe target, ma misurare quasi la stessa informazione.
\end{notaprof}

\section{Wrapper approach}

Il \textbf{wrapper approach} seleziona le feature usando il classificatore come una \emph{black box}. A differenza dei filtri, il wrapper dipende dal classificatore specifico.

Il processo generale è:

\begin{enumerate}
    \item scegliere un subset di feature;
    \item creare training e test set basati su quel subset;
    \item addestrare il classificatore sul training set;
    \item valutare l'accuratezza su un validation set;
    \item ripetere il processo per più subset e scegliere quello con il miglior risultato.
\end{enumerate}

Il vantaggio è che le feature vengono ottimizzate per uno specifico classificatore. Lo svantaggio è il costo computazionale elevato, perché il modello deve essere addestrato e testato molte volte.

\begin{notaprof}
I wrapper sono più lenti dei filtri, ma possono produrre feature subset più adatti a un classificatore specifico.
\end{notaprof}

\section{Search strategies}

Poiché il numero di possibili subset di feature è esponenziale, nei wrapper si usano strategie di ricerca.

\begin{itemize}
    \item \textbf{Exhaustive search}: prova tutte le combinazioni possibili. È completa, ma spesso impraticabile;
    \item \textbf{Forward search}: parte da un insieme vuoto e aggiunge progressivamente la feature che migliora di più il risultato;
    \item \textbf{Backward search}: parte dall'insieme completo e rimuove progressivamente la feature la cui eliminazione danneggia meno l'accuratezza;
    \item \textbf{Greedy search}: prende decisioni locali cercando un buon compromesso tra accuratezza e costo computazionale.
\end{itemize}

\begin{notaslide}
Le slide indicano che la forward search e la backward search hanno costo \(O(n^2 \cdot learning/testing\ time)\) e che esistono anche altre euristiche.
\end{notaslide}

\begin{notaesame}
Se si ha molta fiducia nel set di feature iniziale, può convenire usare \textbf{Backward Search}, perché parte da tutte le feature e ne elimina poche. Se invece si sospetta che molte feature siano inutili o ridondanti, può convenire usare \textbf{Forward Search}, perché costruisce un subset più ristretto partendo da zero.
\end{notaesame}

\section{Train, validation e test}

È fondamentale distinguere tra:

\begin{itemize}
    \item \textbf{training set}: usato per addestrare il modello;
    \item \textbf{validation set}: usato per tuning, scelta delle feature e scelta del modello;
    \item \textbf{test set}: usato solo per la valutazione finale.
\end{itemize}

\begin{notaesame}
Il test set non deve essere usato per scegliere feature o parametri. Se si usa il test set durante la selezione, si introduce data leakage e la valutazione finale non è più realistica.
\end{notaesame}

Il validation set serve a prendere decisioni sul modello senza contaminare il test set. Solo alla fine il modello selezionato viene valutato sul test set.

\section{Feature Selection in Weka}

In Weka, la selezione degli attributi può essere eseguita tramite il pannello \textbf{Select attributes}. I metodi di attribute selection sono composti da due parti:

\begin{itemize}
    \item \textbf{Search Method}: metodo con cui si esplora lo spazio dei subset;
    \item \textbf{Attribute Evaluator}: criterio con cui si valuta la qualità delle feature o dei subset.
\end{itemize}

\begin{notaslide}
Tra i search method citati nelle slide compaiono \textbf{best-first}, \textbf{forward selection}, \textbf{random}, \textbf{exhaustive}, \textbf{genetic algorithm} e \textbf{ranking}.
\end{notaslide}

\begin{notaslide}
Tra gli evaluation method citati nelle slide compaiono metodi \textbf{correlation-based}, \textbf{wrapper}, \textbf{information gain} e \textbf{chi-squared}.
\end{notaslide}

\begin{notaprof}
Weka è flessibile perché permette di combinare quasi arbitrariamente search method ed evaluator.
\end{notaprof}

\section{Lab e collegamento con il progetto}

Il laboratorio richiede di confrontare l'accuratezza con e senza feature selection.

Il workflow richiesto è:

\begin{enumerate}
    \item prendere un dataset;
    \item usare tre classificatori:
    \begin{itemize}
        \item \textbf{RandomForest};
        \item \textbf{NaiveBayes};
        \item \textbf{IBk};
    \end{itemize}
    \item eseguire una \textbf{10-fold cross validation};
    \item confrontare i risultati con e senza feature selection;
    \item discutere se la valutazione è fair;
    \item valutare se usare API o GUI.
\end{enumerate}

\begin{notaesame}
Fare feature selection sull'intero dataset prima della cross-validation non è fair. Significa permettere alla selezione delle feature di vedere anche le fold che dovrebbero essere usate come test. Questo produce data leakage.
\end{notaesame}

\begin{notaprof}
In Weka, per evitare data leakage, è preferibile usare il meta-classificatore \texttt{AttributeSelectedClassifier}. In questo modo la feature selection viene eseguita separatamente dentro ogni fold di training, senza usare la fold di test.
\end{notaprof}

\begin{notaprof}
La GUI è utile per capire il dataset, osservare le distribuzioni e sperimentare rapidamente. Le API o la CLI sono invece più adatte al progetto finale, perché permettono di automatizzare esperimenti ripetibili.
\end{notaprof}

\begin{notaprof}
Se si usano sia feature selection sia balancing, è opportuno prestare attenzione all'ordine delle operazioni. La feature selection dovrebbe essere valutata senza contaminare il test set; anche il balancing deve essere applicato solo sul training set.
\end{notaprof}

\section{Da ricordare}

\begin{notaesame}
In Software Engineering, la feature selection è importante soprattutto perché riduce il costo di misurazione delle metriche, non solo il tempo di training del modello.
\end{notaesame}

\begin{notaesame}
I filtri, come Info Gain e Spearman, sono veloci e indipendenti dal classificatore. I wrapper sono più lenti ma ottimizzano le feature per uno specifico classificatore.
\end{notaesame}

\begin{notaesame}
Il test set non deve mai essere usato per scegliere feature, parametri o modelli. Qualsiasi forma di selezione deve avvenire usando solo training e validation set.
\end{notaesame}

\begin{notaesame}
In Weka, \texttt{AttributeSelectedClassifier} è utile perché consente di eseguire feature selection in modo corretto dentro la cross-validation, evitando data leakage.
\end{notaesame}

\begin{notaesame}
La GUI serve per comprendere e ispezionare il problema; API e CLI servono per automatizzare esperimenti ripetibili e scalabili.
\end{notaesame}
